\documentclass[10pt,a4paper]{article}
\usepackage[margin=1.8cm]{geometry}
\usepackage{booktabs}
\usepackage{array}
\usepackage{xcolor}
\usepackage{colortbl}
\usepackage{caption}
\usepackage{enumitem}
\usepackage{titlesec}
\usepackage{tabularx}
\usepackage{float}
\usepackage{placeins}
\usepackage{hyperref}

\hypersetup{colorlinks=true, linkcolor=blue!60!black, urlcolor=blue!60!black}

\titlespacing*{\section}{0pt}{10pt}{4pt}
\titlespacing*{\subsection}{0pt}{6pt}{2pt}
\setlength{\parskip}{3pt}
\setlength{\parindent}{0pt}

\definecolor{bestrow}{RGB}{220,240,220}
\definecolor{badrow}{RGB}{250,225,225}
\definecolor{neutralrow}{RGB}{240,240,245}
\definecolor{m2row}{RGB}{210,228,255}

\newcommand{\good}[1]{\textcolor{green!50!black}{\textbf{#1}}}
\newcommand{\bad}[1]{\textcolor{red!60!black}{#1}}

\begin{document}

\begin{center}
{\LARGE \textbf{Diagnostic Ablation for M2}}\\[2pt]
{\Large A 2$\times$2 Decomposition of M2's Two Design Choices}\\[3pt]
{\small \textbf{Main method remains M2} (cluster-level ensemble of Direct SSR + LLM-dist, JS = 0.0254). \\
The 2$\times$2 below is an ablation explaining \emph{why} M2 is structured the way it is — not a replacement for it.}\\[2pt]
{\small Compiled 2026-04-28 \ \textbar\ \ Author: mzyy1001 \ \textbar\ \ Model: Qwen3-8B + bge-base-zh-v1.5}\\[2pt]
{\small Sources: \texttt{results/m0\_m2\_pmfs\_6q.json}, \texttt{results/b1\_c2\_pmfs\_6q.json}, \texttt{results/stability\_seeds\_summary.json} (5-seed stability), \texttt{results/flat\_ssr\_6q.json} (flat baseline), \texttt{results/c2\_hardened\_6q.json} (hardened-parser rerun), \texttt{results/topk\_sweep\_6q.json} (top-K sweep)}
\end{center}

\section{Headline (Unchanged): M2 Is the Method}

\begin{table}[H]
\centering
\small
\begin{tabular}{l r r r}
\toprule
\textbf{Method} & \textbf{Mean JS} $\downarrow$ & \textbf{$K_{xy}$} $\uparrow$ & \textbf{$C_{xy}$} $\uparrow$ \\
\midrule
Paper-SSR (Maier transfer, per-consumer flat)  & 0.0430 & 0.827 & 0.877 \\
\textbf{Flat SSR — real posts, no clusters (new control)} & 0.0295 & 0.849 & 0.914 \\
M0 \ Hier cluster-weighted SSR (this work)     & 0.0269 & 0.858 & 0.923 \\
M1 \ Per-cluster LLM-dist (this work)          & 0.0268 & 0.844 & 0.931 \\
\rowcolor{m2row}
\textbf{M2 \ Cluster-level ensemble M0+M1, $w=0.5$ (this work)} & \good{0.0254} & 0.852 & \good{0.932} \\
\bottomrule
\end{tabular}
\caption{M2 remains the headline. The new \emph{Flat SSR} row controls for "what if we use the same real posts but drop the cluster machinery?" — JS = 0.0295. So the gap from Paper-SSR (0.0430) to M2 (0.0254) decomposes as: \textbf{$-31.4\%$ from real posts alone} (Paper $\to$ Flat), \textbf{$-8.8\%$ from cluster-weighting} (Flat $\to$ M0), \textbf{$-5.6\%$ from ensembling LLM-dist} (M0 $\to$ M2). Nothing in this report displaces M2.}
\end{table}

\section{What This Ablation Asks}

M2 fixes \emph{two} design choices simultaneously:

\begin{enumerate}[leftmargin=*,itemsep=0pt]
    \item \textbf{Signal source} — combine the embedding-similarity signal (SSR) with the LLM's direct distribution estimate (LLM-dist), 50/50.
    \item \textbf{Aggregation unit} — for SSR, judge \emph{each post} and average within cluster (M0); for LLM-dist, prompt \emph{the cluster as a whole} and read one PMF (M1).
\end{enumerate}

The 2$\times$2 ablation tests whether either choice could have been made differently. We fill the two missing corners — \textbf{B1} (per-cluster SSR, via mean-pooled centroid) and \textbf{C2} (per-post LLM-dist) — sharing top-10 clusters, $n=50$, seed=42, anchors, and embedder with M0/M1, so all four cells are directly comparable. \emph{All four are diagnostic; M2 remains the headline.}

\section{The 2$\times$2 (Diagnostic Only — None of These Beat M2)}

\begin{table}[H]
\centering
\small
\begin{tabular}{l c c}
\toprule
\textbf{Mean JS over 6Q} & \textbf{Per-cluster} (aggregate $\to$ judge) & \textbf{Per-post} (judge $\to$ aggregate) \\
\midrule
\textbf{SSR (similarity)} & B1 = 0.0280 & \good{M0 = 0.0269} $\leftarrow$ M2 ingredient \\
\textbf{LLM (model PMF)}  & \good{M1 = 0.0268} $\leftarrow$ M2 ingredient & C2(orig) = 0.0259 \  /\  \bad{C2(hard) = 0.0328} \\
\midrule
\multicolumn{3}{l}{\cellcolor{m2row}\textbf{M2 = $0.5 \times $ M0 + $0.5 \times $ M1 (cluster-level ensemble) = 0.0254 — beats every cell here}} \\
\bottomrule
\end{tabular}
\caption{The 2$\times$2 of single-source methods. C2(orig) is the run we originally reported (regex parser, 21\% parse-fail $\to$ uniform PMFs). C2(hard) is the P3 rerun with a multi-strategy parser cascade (10.5\% fail). The hardened number is the fairer single-cell estimate of "per-post LLM" — and it is the worst cell of the 2$\times$2, not competitive with M2.}
\end{table}

\textbf{Why the chosen M2 ingredients (M0 and M1) are reasonable picks for the ensemble:}

\begin{itemize}[leftmargin=*,itemsep=0pt]
    \item \textbf{Within the SSR row, M0 (per-post) is better than B1 (centroid):} $0.0269 < 0.0280$. Mean-pooling embeddings before SSR loses information; per-post SSR keeps it.
    \item \textbf{Within the LLM row, M1 (per-cluster prompt) is more reliable than C2 (per-post prompt):} 0.0268 vs 0.0259 mean is misleading because C2 has two failure modes (see below). On the question where C2 fails (Q4), per-post LLM collapses; per-cluster LLM does not. M1 is the safer ingredient even though C2's mean is nominally lower.
    \item \textbf{Ensembling beats every cell:} M2 = 0.0254 is below all four single cells. The two signal sources (SSR, LLM) carry complementary information — averaging them at cluster level is what gets the headline number.
\end{itemize}

\section{Per-Question — Where M2's Cluster-Level Ensembling Pays Off}

\begin{table}[H]
\centering
\small
\begin{tabular}{l r r r r r r l}
\toprule
\textbf{Question} & \textbf{B1} & \textbf{M0} & \textbf{M1} & \textbf{C2(orig)} & \textbf{C2(hard)} & \textbf{M2} & \textbf{Note} \\
\midrule
Q1 \ 8.0/qsingle\_3   & 0.0148 & 0.0134 & 0.0091 & 0.0012 & 0.0013 & \good{0.0105} & C2 fine here \\
Q2 \ 8.0/qsingle\_4   & 0.0459 & 0.0460 & 0.0506 & 0.0344 & \bad{0.0410} & \good{0.0468} & C2 worse w/o uniform regularizer \\
Q3 \ 9.0/qsingle\_3   & 0.0152 & 0.0198 & \bad{0.0434} & 0.0337 & \bad{0.0418} & \good{0.0304} & C2 hardened $\approx$ M1 (both bad) \\
Q4 \ 9.0/qsingle\_4   & 0.0180 & 0.0173 & 0.0174 & \bad{0.0529} & \bad{0.0531} & \good{0.0158} & C2 collapses regardless of parser \\
Q5 \ 9.0/qsingle\_5   & \bad{0.0538} & 0.0468 & 0.0317 & 0.0263 & \bad{0.0566} & \good{0.0367} & \textbf{Q5 'C2 win' was a parser artifact} \\
Q6 \ 11.0/qsingle\_3  & 0.0205 & 0.0180 & 0.0086 & 0.0070 & \good{0.0029} & \good{0.0125} & Hardened parser actually helps \\
\midrule
\textbf{Mean (6Q)}    & 0.0280 & 0.0269 & 0.0268 & 0.0259 & \bad{0.0328} & \good{0.0254} & M2 wins on the mean \\
\bottomrule
\end{tabular}
\caption{Per-question breakdown. C2(orig) uses the original 21\%-fail parser; C2(hard) uses the new 10.5\%-fail multi-strategy parser. M2 (rightmost) is never the worst on any question, and beats both its parents (M0, M1) on the 6Q mean. The hardened C2 numbers reveal that the original C2 mean was inflated by parse-fail-induced uniform-PMF fallback on Q2/Q3/Q5; once the parser is fixed, C2 is the worst single-source baseline of all.}
\end{table}

\textbf{What the per-question table shows about M2's design:}

\begin{itemize}[leftmargin=*,itemsep=0pt]
    \item \textbf{Q4 (dampness season) — vindicates per-cluster LLM (M1) over per-post LLM (C2).} C2 collapses to 54.8\% mass on the "summer" option (JS=0.0529); M1 stays calibrated (JS=0.0174). M2 inherits M1 here and lands at \good{0.0158} — best of any method on Q4. If C2 had been the LLM ingredient, M2 would have been worse here.
    \item \textbf{Q3 (dampness cause) — vindicates ensembling.} M1 alone is bad (0.0434); M0 is fine (0.0198); M2 partially repairs the LLM error by averaging (0.0304). The complementarity argument for M2 is most visible here.
    \item \textbf{Q1, Q5, Q6 — M2 is not the per-question optimum but is robust.} Different single cells win on different questions (C2 on Q1/Q5/Q6, M0 on Q4, B1 on Q3) — no single cell dominates. M2 averages these heterogeneous patterns into a uniformly competent number.
\end{itemize}

\section{C2 Reliability Audit + Hardened-Parser Rerun (P3)}

\subsection{The original C2 parser was leaking}

C2 generates one JSON PMF per post; failed parses fall back to a uniform PMF.

\begin{table}[H]
\centering
\small
\begin{tabular}{l r r r r r}
\toprule
\textbf{Question} & \textbf{Original fails} & \textbf{Hardened fails} & \textbf{Original JS} & \textbf{Hardened JS} & \textbf{$\Delta$ JS} \\
\midrule
Q1 \ 8.0/qsingle\_3   & 18.2\% & \good{10.8\%} & 0.0012 & 0.0013 & $+0.0001$ \\
Q2 \ 8.0/qsingle\_4   & 24.8\% & \good{15.2\%} & 0.0344 & 0.0410 & \bad{$+0.0066$} \\
Q3 \ 9.0/qsingle\_3   & 10.2\% & \good{4.8\%}  & 0.0337 & 0.0418 & \bad{$+0.0082$} \\
Q4 \ 9.0/qsingle\_4   & 3.4\%  & \good{2.2\%}  & 0.0529 & 0.0531 & $+0.0002$ \\
\rowcolor{badrow}
Q5 \ 9.0/qsingle\_5   & \bad{53.2\%} & \good{22.2\%} & 0.0263 & 0.0566 & \bad{$+0.0303$} \\
Q6 \ 11.0/qsingle\_3  & 15.8\% & \good{8.0\%}  & 0.0070 & 0.0029 & \good{$-0.0041$} \\
\midrule
\textbf{Mean (6Q)}    & 20.9\% & \good{10.5\%} & 0.0259 & \bad{\textbf{0.0328}} & \bad{\textbf{$+0.0069$}} \\
\bottomrule
\end{tabular}
\caption{Hardened parser cuts the fail rate roughly in half on every question, but C2's mean JS gets \emph{worse} (0.0259 $\to$ 0.0328). The original C2 mean was bolstered by parse-fail-induced uniform-PMF fallback, which acted as an inadvertent regularizer on questions where the LLM's actual per-post answers were confidently wrong.}
\end{table}

\subsection{Strategy hits across 3000 LLM calls}

\begin{table}[H]
\centering
\small
\begin{tabular}{l r r}
\toprule
\textbf{Parser strategy} & \textbf{Hits} & \textbf{Share} \\
\midrule
S1 \ Full JSON object \texttt{\{"distribution":[\ldots]\}} & 2412 & 80.4\% \\
S2 \ Bare \texttt{[a,b,c,\ldots]} with exact $n_{\rm opts}$ entries  & 3 & 0.1\% \\
S3 \ Last $n_{\rm opts}$ numeric tokens anywhere & 269 & 9.0\% \\
S4 \ Percent-style lines (e.g. ``option1: 30\%''\ldots) & 0 & 0\% \\
\midrule
\rowcolor{badrow}
\textbf{Fail (uniform PMF fallback)} & 316 & 10.5\% \\
\bottomrule
\end{tabular}
\caption{The original parser only used a regex roughly equivalent to S2, hitting 0.1\% of calls. S1 (proper JSON-object extraction) carries 80\% of the load; S3 (last-$n$-nums) catches another 9\% where the LLM omitted the brackets. S4 was unused. The 10.5\% remaining failures are responses with neither JSON nor parseable numerics — those are genuine LLM rejections (prompt refusal, off-topic answers) that uniform fallback cannot rescue.}
\end{table}

\subsection{What this means for the report}

\begin{itemize}[leftmargin=*,itemsep=1pt]
    \item \textbf{The original C2 result (0.0259) was misleadingly low.} It was within striking distance of M2 (0.0254), but only because $\sim$21\% of calls returned uniform PMFs that happened to dampen the LLM's confidently-wrong answers on Q2/Q3/Q5.
    \item \textbf{C2 hardened (0.0328) is no longer competitive with M2 (0.0254), or even with M0 (0.0269) or Flat SSR (0.0295).} It is also worse than the Paper-SSR transfer baseline on Q5 specifically.
    \item \textbf{Q4's mode collapse is real, not a parser artifact} — fail rate was already only 3\% in the original; hardening barely moves JS (0.0529 $\to$ 0.0531).
    \item \textbf{Q6 is the one question where the hardened parser improves C2} — fewer fails capture more of a genuinely strong per-post LLM signal here.
    \item \textbf{This strengthens M2's design choice} of M1 (per-cluster prompting) as the LLM ingredient. Per-post LLM is unreliable in two compounding ways: confidence collapses on some Qs (Q2, Q3, Q4) \emph{and} JSON parsing is fragile (Q5). Per-cluster prompting avoids both.
\end{itemize}

\section{Stability of M0 / M1 / M2 Across 5 Random Seeds}

We re-ran the M0 + M1 + M2 pipeline with five distinct seeds (42, 123, 456, 789, 1024), varying both post-sampling and torch RNG; anchors stayed deterministic across runs to avoid confounding.

\begin{table}[H]
\centering
\small
\begin{tabular}{l c c c c c c}
\toprule
\textbf{Method} & \textbf{Mean} & \textbf{Std} & \textbf{Min} & \textbf{Max} & \textbf{Range} & \textbf{Original (seed=42 dump)} \\
\midrule
M0 hier cluster-weighted SSR & 0.0271 & \good{0.0002} & 0.0269 & 0.0275 & 0.0006 & 0.0269 \\
M1 per-cluster LLM-dist      & 0.0260 & 0.0011 & 0.0243 & 0.0271 & 0.0028 & 0.0268 \\
\rowcolor{m2row}
\textbf{M2 ensemble (M0 + M1, w=0.5)} & \good{0.0252} & \good{0.0004} & 0.0245 & 0.0256 & 0.0011 & 0.0254 \\
\bottomrule
\end{tabular}
\caption{Mean JS over 6Q across 5 seeds. M0 is essentially deterministic ($\sigma$ = 0.0002, post-sampling only). M1 is the noisiest ingredient because of LLM stochasticity ($\sigma$ = 0.0011). M2 inherits low variance ($\sigma$ = 0.0004) because ensembling damps M1's noise. \textbf{The M2-vs-M0 gap (mean 0.0019, $\approx$ 4.7$\sigma$ of M2's std) is robust across seeds.}}
\end{table}

\textbf{Per-question M2 stability (mean $\pm$ std over 5 seeds):}

\begin{table}[H]
\centering
\small
\begin{tabular}{l c c c c c c}
\toprule
& \textbf{Q1} & \textbf{Q2} & \textbf{Q3} & \textbf{Q4} & \textbf{Q5} & \textbf{Q6} \\
\textbf{M2 JS} & $0.008 \pm 0.002$ & $0.045 \pm 0.001$ & $0.032 \pm 0.002$ & $0.015 \pm 0.001$ & $0.037 \pm 0.002$ & $0.013 \pm 0.001$ \\
\bottomrule
\end{tabular}
\caption{Per-question M2 JS stability. Std is $\le$ 0.0020 on every question. The Q3 ensemble compromise (M2 worsens Q3 vs M0 by $\sim$50\%) is consistently reproducible across seeds — it is a true property of the ensemble, not a single-seed artifact.}
\end{table}

\section{Top-K Sweep — Is K=10 a Good Choice?}

We re-ran the M0+M1+M2 pipeline with $K \in \{5, 10, 15, 20\}$ (seed 42, anchors fixed across K). The chosen $K=10$ default sits at or one step below the optimum.

\begin{table}[H]
\centering
\small
\begin{tabular}{c r r r r}
\toprule
\textbf{K} & \textbf{M0} JS & \textbf{M1} JS & \textbf{M2} JS & \textbf{$\Delta$ M2 vs K=10} \\
\midrule
5  & 0.0270 & \bad{0.0306} & 0.0269 & $+0.0019$ ($+7.6\%$) \\
\rowcolor{m2row}
\textbf{10 (default)} & 0.0269 & 0.0275 & \good{0.0250} & --- \\
15 & 0.0275 & \good{0.0250} & \good{0.0249} & $-0.0001$ (tie) \\
20 & 0.0277 & 0.0265 & 0.0258 & $+0.0008$ ($+3.2\%$) \\
\bottomrule
\end{tabular}
\caption{M2 top-K sweep. M2 is bowl-shaped: $K=10$ and $K=15$ tie at 0.0250 / 0.0249; $K=5$ is materially worse (too few clusters); $K=20$ regresses (cluster-mass weights spread too thin onto noisy tail clusters). \textbf{The chosen $K=10$ is within 0.0001 of optimal} — $K=15$ does not deliver a meaningful improvement and is more expensive ($+50\%$ LLM calls).}
\end{table}

\textbf{Per-question M2 across K} — where the K=10 vs K=15 tie splits unevenly:

\begin{table}[H]
\centering
\small
\begin{tabular}{l r r r r l}
\toprule
\textbf{Question} & \textbf{K=5} & \textbf{K=10} & \textbf{K=15} & \textbf{K=20} & \textbf{Best K} \\
\midrule
Q1 \ 8.0/qsingle\_3   & 0.0115 & \good{0.0066} & 0.0068 & 0.0113 & K=10 \\
Q2 \ 8.0/qsingle\_4   & 0.0442 & 0.0462 & \good{0.0436} & 0.0450 & K=15 \\
Q3 \ 9.0/qsingle\_3   & 0.0355 & 0.0319 & \good{0.0298} & 0.0302 & K=15 \\
Q4 \ 9.0/qsingle\_4   & 0.0165 & 0.0192 & 0.0160 & \good{0.0156} & K=20 \\
Q5 \ 9.0/qsingle\_5   & 0.0406 & \good{0.0328} & \bad{0.0418} & 0.0393 & K=10 \\
Q6 \ 11.0/qsingle\_3  & 0.0131 & 0.0135 & \good{0.0117} & 0.0134 & K=15 \\
\midrule
\textbf{Mean (6Q)}    & 0.0269 & \good{0.0250} & \good{0.0249} & 0.0258 & --- \\
\bottomrule
\end{tabular}
\caption{Per-question M2 JS across K. K=15 wins 4 of 6 questions (Q2, Q3, Q4, Q6) and K=10 wins 2 (Q1, Q5). However, Q5 has a sharp regression at K=15 ($0.0328 \to 0.0418$, $+27\%$) that almost exactly cancels K=15's wins on the other five — explaining why the means are tied. The K=10 default is robust because it never has the K=15 Q5 failure mode.}
\end{table}

\textbf{Read:} $K=10$ is a sweet spot, not a knob worth re-tuning. The bowl shape (worst at $K=5$, second-worst at $K=20$, two-way tie at $K=10/15$) suggests the cluster-mass weighting itself is doing meaningful work — extending top-K dilutes per-cluster signal-to-noise without commensurate coverage gain.

\section{Summary}

\begin{itemize}[leftmargin=*,itemsep=2pt]
    \item \textbf{M2 (cluster-level ensemble of Direct SSR + LLM-dist, $w=0.5$) is the headline method, JS = 0.0254 ($\pm$ 0.0004 over 5 seeds).} Unchanged.
    \item \textbf{Flat SSR control (new):} the gap from Paper-SSR (0.0430) to M2 (0.0254) decomposes — real posts alone get to 0.0295 ($-31\%$); cluster-weighting takes it to 0.0269 ($-9\%$); LLM-dist ensembling takes it to 0.0254 ($-6\%$). The biggest single lever is using real posts; clustering and ensembling each add a smaller, consistent gain.
    \item \textbf{5-seed stability (new):} M0 is essentially deterministic ($\sigma$ = 0.0002), M2 is also stable ($\sigma$ = 0.0004), and the M2 $-$ M0 gap is robust ($\approx$ 4.7$\sigma$).
    \item The 2$\times$2 ablation explains \emph{why} M2's two ingredients (M0 = per-post SSR, M1 = per-cluster LLM-dist) are the right choices: B1 (centroid SSR) loses information vs M0; C2 (per-post LLM-dist) is unstable on Q4/Q5 vs M1.
    \item No cell of the 2$\times$2 — including the nominally-best C2 — beats M2 on the 6Q mean.
    \item B1 and C2 are reported as \emph{ablation evidence supporting M2's design}, not as candidate replacements for M2.
    \item \textbf{P3 C2-hardened-parser rerun (new):} multi-strategy JSON parser cascade reduces fail rate from 21\% to 10.5\%. C2's mean JS jumps from 0.0259 to \textbf{0.0328} — i.e. the original C2 number was inflated by parse-fail $\to$ uniform-PMF fallback that happened to dampen the LLM's confidently-wrong per-post answers on Q2/Q3/Q5. Q5 is the most striking case: hardened JS = 0.0566 vs original 0.0263 (the "C2 wins on Q5" claim does not survive). With proper parsing, no cell of the 2$\times$2 is competitive with M2.
    \item \textbf{Top-K sweep (new):} M2 is bowl-shaped over $K \in \{5, 10, 15, 20\}$ with means 0.0269 / 0.0250 / 0.0249 / 0.0258. The chosen $K=10$ is within 0.0001 of optimal and avoids a Q5-specific regression that K=15 exhibits ($+27\%$ on Q5 alone). $K=10$ stays as the default; not a tunable.
\end{itemize}

\end{document}
